\documentclass[a4paper,11pt]{article}
\usepackage[utf8]{inputenc}
\usepackage[dutch]{babel}
\usepackage{array}
\usepackage[margin=2cm]{geometry}
\pagestyle{empty}

\usepackage{fancyhdr}  % for footer, header
\pagestyle{fancy}
\fancyhf{} % clear all header and footer fields
\renewcommand{\headrulewidth}{0pt}  % clear hrule in header
\cfoot{\fontsize{6pt}{7.2pt}\selectfont autogenerated by nwc-concat \hspace{0.5cm} \today }
\chead{\fontsize{8pt}{9.6pt}\selectfont Returnability (333) structuur.pdf}

\begin{document}

\section*{Lied structuur}

\begin{tabular}{ll}
\hline
\textbf{Basis} & \\
\hline
Titel & Returnability (333) \\
Maatsoort & 4/4 \\
Tempo & 184 \\
\#Maten & 173 \\
Duur & 3:46 \\
\hline
\end{tabular}

\vspace{0.5cm}

\subsection*{Lied delen}

\begin{tabular}{l|c|p{8cm}}
\hline
\textbf{Naam} & \textbf{\#Maten} & \textbf{Akkoorden (\#mt)} \\
\hline
intro & 8 & A7sus2(8) \\
couplet 1 & 18 & A7sus2(5), Dmaj7(4), A7sus2(5), Dmaj7(4) \\
refrein & 15 & Bm6(3), Em(4), Edim(add 11)(2), Em (add9+11)(2), Dmaj7(2), B(2) \\
overgang refr-couplet & 10 & E6(2), A (schuif)(4), A7sus2(4) \\
couplet & 24 & A7sus2(4), Dmaj7(4), A7sus2(4), Dmaj7(4), A7sus2(4), Dmaj7(4) \\
middenstuk & 36 & Em7(4), A7(4), Em7(4), A7(4), Em7(4), A7(4), Em7(4), A7(4), A7sus2(4) \\
uittro & 8 & Em6(2), A(4), D(2) \\
\hline
\end{tabular}

\subsection*{Compositie}

Compositie van het lied, met tussen haakjes het aantal maten:

\vspace{0.3cm}

\renewcommand{\arraystretch}{1.3}  % some space between rows

\begin{tabular}{l|l|c|p{8cm}}
\hline
\textbf{Volgnr} & \textbf{Deel} & \textbf{\#Maten} & \textbf{Akkoorden(\#mt)} \\
\hline
1 & intro & 8 & A7sus2(8) \\
2 & couplet 1 & 18 & A7sus2(5), Dmaj7(4), A7sus2(5), Dmaj7(4) \\
3 & refrein & 15 & Bm6(3), Em(4), Edim(add 11)(2), Em (add9+11)(2), Dmaj7(2), B(2) \\
4 & overgang refr-couplet & 10 & E6(2), A (schuif)(4), A7sus2(4) \\
5 & couplet & 24 & A7sus2(4), Dmaj7(4), A7sus2(4), Dmaj7(4), A7sus2(4), Dmaj7(4) \\
6 & refrein & 15 & Bm6(3), Em(4), Edim(add 11)(2), Em (add9+11)(2), Dmaj7(2), B(2) \\
7 & middenstuk & 36 & Em7(4), A7(4), Em7(4), A7(4), Em7(4), A7(4), Em7(4), A7(4), A7sus2(4) \\
8 & couplet & 24 & A7sus2(4), Dmaj7(4), A7sus2(4), Dmaj7(4), A7sus2(4), Dmaj7(4) \\
9 & refrein & 15 & Bm6(3), Em(4), Edim(add 11)(2), Em (add9+11)(2), Dmaj7(2), B(2) \\
10 & uittro & 8 & Em6(2), A(4), D(2) \\
\hline
\end{tabular}

\end{document}
